\documentclass[10pt]{article}

\usepackage[utf8]{inputenc}
\usepackage[spanish]{babel}
\usepackage{amsmath}
\usepackage[margin=2.5cm]{geometry}
\usepackage{graphicx}
\usepackage{tikz}
\usetikzlibrary{babel}
\usepackage{array}
\usepackage[table]{xcolor}
\usepackage{float}
\usepackage{fancyhdr}
\usepackage[colorlinks=true,linkcolor=black,urlcolor=black,citecolor=black]{hyperref}

\definecolor{tablebackground}{rgb}{0.8, 0, 0}
\newcolumntype{x}[1]{>{\centering\arraybackslash\hspace{0pt}}p{#1}}
\newcolumntype{M}[1]{>{\centering\arraybackslash}m{#1}}
\newcolumntype{N}{@{}m{0pt}@{}}

% Datos de portada y encabezado
\newcommand{\itemEmail}{ccapatinta@unsa.edu.pe}
\newcommand{\itemStudent}{Camilo B.A. Capatinta Almanza (Teoría grupo B)}
\newcommand{\itemCourse}{Arq. de computadoras}
\newcommand{\itemCourseCode}{1702227}
\newcommand{\itemSemester}{IV}
\newcommand{\itemUniversity}{Universidad Nacional de San Agustín de Arequipa}
\newcommand{\itemFaculty}{Facultad de Ingeniería de Producción y Servicios}
\newcommand{\itemDepartment}{Departamento Académico de Ingeniería de Sistemas e Informática}
\newcommand{\itemSchool}{Escuela Profesional de Ingeniería de Sistemas}
\newcommand{\itemAcademic}{2026 - B}
\newcommand{\itemInput}{08/11/2026}
\newcommand{\itemOutput}{08/11/2026}
\newcommand{\itemPracticeNumber}{02}
\newcommand{\itemTheme}{PARTE II}

\pagestyle{fancy}
\fancyhf{}
\setlength{\headheight}{40.51407pt}
\renewcommand{\headrulewidth}{1pt}
\renewcommand{\footrulewidth}{1pt}

\fancyhead[L]{\raisebox{-0.2\height}{\includegraphics[width=3cm]{img/logo_episunsa.png}}}
\fancyhead[C]{\fontsize{7}{7}\selectfont    \itemUniversity \\ \itemFaculty \\ \itemDepartment \\ \itemSchool \\ \textbf{\itemCourse}}
\fancyhead[R]{\raisebox{-0.2\height}{\includegraphics[width=1.2cm]{img/logo_abet.png}}}

\fancyfoot[L]{Estudiante Camilo Capatinta}
\fancyfoot[C]{\itemCourse}
\fancyfoot[R]{Página \thepage}

\begin{document}

	\vspace*{10px}
	
	\begin{center}
		\parbox{0.8\textwidth}{\centering\fontsize{17}{20}\selectfont\textbf{EXAMEN PARCIAL \itemPracticeNumber}}
	\end{center}
	\begin{center}
		\parbox{0.8\textwidth}{\centering\textbf{\Large Tema: \itemTheme}}
	\end{center}

\begin{flushright}
    \begin{tabular}{|M{2.5cm}|N|}
        \hline 
        \rowcolor{tablebackground}
        \color{white} \textbf{Nota}  \\
        \hline 
             \\[30pt]
        \hline          
    \end{tabular}
\end{flushright}    

\begin{table}[h!]
    \begin{tabular}{|x{4.7cm}|x{4.8cm}|x{4.8cm}|}        
        \hline 
        \rowcolor{tablebackground}
        \color{white} \textbf{Estudiante} & \color{white}\textbf{Escuela}  & \color{white}\textbf{Asignatura}   \\
        \hline 
        {\itemStudent \par \itemEmail} & \itemSchool & {\itemCourse \par Semestre: \itemSemester \par Código: \itemCourseCode}     \\
        \hline          
    \end{tabular}
\end{table}     

\begin{table}[H]
    \begin{tabular}{|x{4.7cm}|x{4.8cm}|x{4.8cm}|}
        \hline 
        \rowcolor{tablebackground}
        \color{white}\textbf{Laboratorio} & \color{white}\textbf{Tema}  & \color{white}\textbf{Duración}   \\
        \hline 
        \itemPracticeNumber & \itemTheme & 02 horas   \\
        \hline 
    \end{tabular}
\end{table}

\begin{table}[H]
    \begin{tabular}{|x{4.7cm}|x{4.8cm}|x{4.8cm}|}
        \hline 
        \rowcolor{tablebackground}
        \color{white}\textbf{Semestre académico} & \color{white}\textbf{Fecha de inicio}  & \color{white}\textbf{Fecha de entrega}   \\
        \hline 
        \itemAcademic & \itemInput &  \itemOutput  \\
        \hline 
    \end{tabular}
\end{table}

\newpage
\tableofcontents
\thispagestyle{plain}
\clearpage

\section{Pregunta 1: Decimal a IEEE 754 de precisión simple}
 
Se representa $13.15625_{10}$ en el formato IEEE 754 de precisión simple (32 bits: 1 bit de signo, 8 bits de exponente con sesgo 127 y 23 bits de mantisa).
 
\begin{itemize}
  \item \textbf{Signo:} el número es positivo, por lo tanto $S = 0$.
  \item \textbf{Parte entera:} $13_{10} = 1101_2$.
  \item \textbf{Parte fraccionaria:} se multiplica por 2 de forma sucesiva y se toma la parte entera de cada resultado.
\end{itemize}
 
\begin{center}
\begin{tabular}{c c c}
  \hline
  \textbf{Operación} & \textbf{Resultado} & \textbf{Bit} \\
  \hline
  $0.15625 \times 2$ & $0.3125$ & 0 \\
  $0.3125 \times 2$  & $0.625$  & 0 \\
  $0.625 \times 2$   & $1.25$   & 1 \\
  $0.25 \times 2$    & $0.5$    & 0 \\
  $0.5 \times 2$     & $1.0$    & 1 \\
  \hline
\end{tabular}
\end{center}
 
\begin{itemize}
  \item Entonces $0.15625_{10} = 0.00101_2$ y $13.15625_{10} = 1101.00101_2$.
  \item \textbf{Normalización:} $1101.00101_2 = 1.10100101_2 \times 2^{3}$ (la coma se desplaza 3 lugares a la izquierda), es decir, $e = 3$.
  \item \textbf{Exponente sesgado:} $E = e + 127 = 3 + 127 = 130_{10} = 10000010_2$.
  \item \textbf{Mantisa:} son los bits que siguen al 1 implícito, $10100101_2$, completados con ceros hasta los 23 bits.
\end{itemize}
 
\begin{center}
\begin{tabular}{|c|c|c|}
  \hline
  \textbf{Signo (1 bit)} & \textbf{Exponente (8 bits)} & \textbf{Mantisa (23 bits)} \\
  \hline
  0 & 10000010 & 10100101000000000000000 \\
  \hline
\end{tabular}
\end{center}
 
\textbf{Respuesta:} $13.15625_{10} = 0\ 10000010\ 10100101000000000000000_2 = 41528000_{16}$.
 
\textbf{Comprobación:} $1.10100101_2 = 1 + \tfrac{1}{2} + \tfrac{1}{8} + \tfrac{1}{64} + \tfrac{1}{256} = 1.64453125$ y $1.64453125 \times 2^{3} = 13.15625$.
 
% ----------------------------------------------------------------------
\section{Pregunta 2: IEEE 754 de precisión simple a decimal}
 
Número dado: $1\ 01100110\ 10100000000000000000000$ (signo, exponente y mantisa).
 
\begin{itemize}
  \item \textbf{Signo:} $S = 1$, por lo tanto el número es negativo.
  \item \textbf{Exponente:} $E = 01100110_2 = 64 + 32 + 4 + 2 = 102_{10}$, entonces $e = E - 127 = 102 - 127 = -25$.
  \item \textbf{Mantisa:} $M = 101\,000\ldots0_2$, por lo que el significando (con el 1 implícito) es $1.101_2 = 1 + \tfrac{1}{2} + \tfrac{1}{8} = 1.625$.
\end{itemize}
 
\[
  N = (-1)^{S} \times 1.M \times 2^{e} = -1.625 \times 2^{-25} = -\frac{13}{2^{28}} = -\frac{13}{268\,435\,456}
\]
 
\textbf{Respuesta:} $N \approx -4.8429 \times 10^{-8}$ (en hexadecimal, $B3500000_{16}$).
 
% ----------------------------------------------------------------------
\section{Pregunta 3: Memoria jerárquica}
 
La \textit{memoria jerárquica} organiza la memoria del computador en varios niveles construidos con tecnologías distintas. Los niveles cercanos al procesador son pequeños, muy rápidos y caros por bit; los niveles lejanos son grandes, lentos y baratos. Como ninguna tecnología es a la vez rápida, barata y de gran capacidad, se combinan para que el sistema se comporte como una memoria grande y casi tan rápida como la más veloz.
 
La jerarquía funciona gracias al \textit{principio de localidad}:
 
\begin{itemize}
  \item \textbf{Localidad temporal:} un dato usado recientemente tiende a volver a usarse pronto.
  \item \textbf{Localidad espacial:} después de acceder a una dirección es probable acceder a las cercanas, por eso se copian bloques completos y no palabras sueltas.
\end{itemize}
 
Cuando el procesador necesita un dato, lo busca primero en el nivel más cercano. Si está (acierto) lo usa; si no está (fallo), se trae un bloque desde el nivel siguiente. Cada nivel contiene un subconjunto de los datos del nivel inferior.
 
\begin{center}
\begin{minipage}{\linewidth}\centering
\begin{tikzpicture}[
    nivel/.style={draw, rectangle, minimum width=6.4cm, minimum height=1cm,
                  align=center, font=\small},
    enlace/.style={<->, thick}
  ]
  \node[nivel] (cpu)   at (0,0)    {\textbf{Procesador (CPU)}\\ registros};
  \node[nivel] (cache) at (0,-1.9) {\textbf{Memoria caché} (L1, L2, L3)\\ SRAM};
  \node[nivel] (ram)   at (0,-3.8) {\textbf{Memoria principal} (RAM)\\ DRAM};
  \node[nivel] (sec)   at (0,-5.7) {\textbf{Almacenamiento secundario}\\ SSD / HDD};
  \node[nivel] (ter)   at (0,-7.6) {\textbf{Almacenamiento terciario}\\ cintas, discos ópticos, nube};
 
  \draw[enlace] (cpu)   -- node[right, font=\scriptsize] {palabras} (cache);
  \draw[enlace] (cache) -- node[right, font=\scriptsize] {bloques}  (ram);
  \draw[enlace] (ram)   -- node[right, font=\scriptsize] {páginas}  (sec);
  \draw[enlace] (sec)   -- node[right, font=\scriptsize] {archivos} (ter);
 
  \draw[->, thick] (-4.1,-7.6) -- (-4.1,0)
        node[midway, above, sloped, font=\scriptsize] {más rápida, más cara, menor capacidad};
  \draw[->, thick] (4.1,0) -- (4.1,-7.6)
        node[midway, above, sloped, font=\scriptsize] {más lenta, más barata, mayor capacidad};
\end{tikzpicture}
 
\small Esquema de un procesador conectado a una memoria jerárquica.
\end{minipage}
\end{center}
 
Niveles que se encuentran:
 
\begin{itemize}
  \item \textbf{Registros de la CPU:} están dentro del procesador y guardan operandos y resultados inmediatos. Son el almacenamiento más rápido, pero muy pocos.
  \item \textbf{Memoria caché (L1, L2, L3):} memoria SRAM entre la CPU y la RAM que guarda copias de los bloques más usados. La administra el hardware; L1 es la más pequeña y rápida, y L3 la más grande y lenta.
  \item \textbf{Memoria principal (RAM):} memoria DRAM que contiene los programas y datos en ejecución. El sistema operativo la administra con memoria virtual (paginación y MMU).
  \item \textbf{Secundario (SSD o HDD):} memoria no volátil que guarda archivos y programas y recibe las páginas que no caben en la RAM.
  \item \textbf{Almacenamiento terciario:} cintas, discos ópticos o la nube, para copias de seguridad y datos que casi no se usan.
\end{itemize}
 
\begingroup
\footnotesize
\setlength{\tabcolsep}{4pt}
\begin{center}
\begin{tabular}{|l|l|l|l|l|}
  \hline
  \textbf{Nivel} & \textbf{Tecnología} & \textbf{Acceso} & \textbf{Capacidad} & \textbf{Gestión} \\
  \hline
  Registros & Flip-flops & $<$ 1 ns & cientos de bytes & Compilador \\
  \hline
  Caché (L1--L3) & SRAM & 1 a 30 ns & KB a MB & Hardware \\
  \hline
  Memoria principal & DRAM & 50 a 100 ns & GB & S.O. \\
  \hline
  Secundario & SSD / HDD & 0.1 a 10 ms & GB a TB & S.O. \\
  \hline
  Terciario & Cinta / nube & s a min & TB a PB & Operador \\
  \hline
\end{tabular}
\end{center}
Valores típicos y orientativos (S.O. = sistema operativo).
\endgroup
 
Al bajar en la jerarquía aumentan la capacidad y el tiempo de acceso, y disminuye el costo por bit.
 
% ----------------------------------------------------------------------
\section{Pregunta 4: Ejecución de la instrucción \texttt{SW \#2AB}}
 
La instrucción \texttt{SW \#2AB} (\emph{store word}) guarda el contenido del acumulador en la posición de memoria $2AB_{16}$ (el símbolo \# indica que $2AB$ es una dirección de memoria). En notación de transferencia de registros:
\[
  M[2AB] \leftarrow AC
\]
 
Según la secuencia del computador elemental explicada en clase, la instrucción se ejecuta en cuatro fases: búsqueda (\emph{fetch}), decodificación y lectura de operandos, ejecución de la operación y almacenamiento del resultado. Como la instrucción tiene 16 bits y la memoria transfiere 8 bits por acceso, la búsqueda requiere dos lecturas consecutivas. Los elementos que intervienen son:
 
\begin{itemize}
  \item \textbf{PC (contador de programa):} contiene la dirección de la siguiente instrucción.
  \item \textbf{MAR (registro de dirección de memoria):} coloca la dirección en el bus de direcciones.
  \item \textbf{MDR (registro de datos de memoria):} recibe o entrega el dato que viaja por el bus de datos.
  \item \textbf{IR (registro de instrucción):} guarda la instrucción que se está ejecutando.
  \item \textbf{AC (acumulador):} contiene el dato que se va a almacenar.
  \item \textbf{Unidad de control (UC):} decodifica la instrucción y genera las señales de control.
\end{itemize}
 
\begin{center}
\begin{minipage}{\linewidth}\centering
\begin{tikzpicture}[
    >=stealth,
    reg/.style={draw, rectangle, minimum width=1.8cm, minimum height=0.8cm, font=\small},
    etq/.style={font=\scriptsize}
  ]
  \node[reg] (pc)  at (0,0)      {PC};
  \node[reg] (mar) at (3.4,0)    {MAR};
  \node[reg] (ir)  at (0,-2.6)   {IR};
  \node[reg] (mdr) at (3.4,-2.6) {MDR};
  \node[reg] (ac)  at (3.4,-5)   {AC};
  \node[reg, align=center] (uc) at (0,-5) {Unidad de\\control};
  \node[draw, rectangle, minimum width=2.2cm, minimum height=3.6cm] (mem) at (7.6,-1.3) {};
  \node[font=\footnotesize] at (7.6,0.1) {Memoria $M$};
  \node[draw, rectangle, minimum width=1.6cm, font=\scriptsize] at (7.6,-1.3) {$M[2AB]$};
 
  \draw[->]  (pc)  -- node[above, etq] {dirección} (mar);
  \draw[->]  (mar.east) -- node[above, etq] {dirección} (mem.west |- mar.east);
  \draw[<->] (mdr.east) -- node[above, etq] {lectura / escritura} (mem.west |- mdr.east);
  \draw[->]  (mdr) -- node[above, etq] {byte} (ir);
  \draw[->]  (ir.north east) -- node[midway, right, etq] {campo dir.} (mar.south west);
  \draw[->]  (ac)  -- node[right, etq] {dato} (mdr);
  \draw[->]  (ir)  -- node[right, etq] {decodifica} (uc);
\end{tikzpicture}
 
\small Registros y transferencias que intervienen en el ciclo de instrucción.
\end{minipage}
\end{center}
 
\begingroup
\small
\begin{center}
\begin{tabular}{|c|l|p{6.2cm}|}
  \hline
  \textbf{Paso} & \textbf{Micro-operación} & \textbf{Qué ocurre} \\
  \hline
  \multicolumn{3}{|l|}{\textbf{1. Búsqueda de la instrucción (\emph{fetch})}} \\
  \hline
  1 & $MAR \leftarrow PC$ & Se coloca en MAR la dirección del primer byte de la instrucción. \\
  \hline
  2 & $MDR \leftarrow M[MAR]$ & La memoria entrega el primer byte (8 bits) a MDR. \\
  \hline
  3 & $IR[\text{byte 1}] \leftarrow MDR$ & El primer byte se guarda en la parte correspondiente de IR. \\
  \hline
  4 & $PC \leftarrow PC + 1$ & PC apunta al segundo byte de la instrucción. \\
  \hline
  5 & $MAR \leftarrow PC$ & Se coloca en MAR la dirección del segundo byte. \\
  \hline
  6 & $MDR \leftarrow M[MAR]$ & La memoria entrega el segundo byte (8 bits) a MDR. \\
  \hline
  7 & $IR[\text{byte 2}] \leftarrow MDR$ & Se completa la instrucción de 16 bits en IR. \\
  \hline
  8 & $PC \leftarrow PC + 1$ & PC apunta a la instrucción siguiente. \\
  \hline
  \multicolumn{3}{|l|}{\textbf{2. Decodificación y lectura de operandos}} \\
  \hline
  -- & Decodificar IR & La UC reconoce \texttt{SW} y obtiene la dirección $2AB$ del campo de la instrucción. El dato fuente ya está en AC. \\
  \hline
  \multicolumn{3}{|l|}{\textbf{3. Ejecución de la operación}} \\
  \hline
  9 & $MAR \leftarrow IR[\text{dir}]$ & La dirección de destino ($2AB$) se coloca en MAR. \\
  \hline
  10 & $MDR \leftarrow AC$ & El dato (o el byte que corresponda) se coloca en MDR para transferirlo a memoria. \\
  \hline
  \multicolumn{3}{|l|}{\textbf{4. Almacenamiento del resultado}} \\
  \hline
  11 & $M[MAR] \leftarrow MDR$ & Se activa la escritura y se almacena el dato en la dirección indicada. Si la palabra ocupa más de 8 bits, se repite la transferencia para sus bytes restantes. \\
  \hline
\end{tabular}
\end{center}
\endgroup
 
La búsqueda realiza dos lecturas de memoria porque cada acceso entrega 8 bits y la instrucción ocupa 16 bits; por ello PC se incrementa dos veces. La dirección $2AB$ se obtiene al decodificar el campo de dirección y el dato fuente se toma de AC. En la fase final se escribe en memoria; si el tamaño del dato almacenado supera los 8 bits, se necesitan transferencias adicionales para los bytes restantes, de acuerdo con el formato y el orden de bytes definidos para el computador. El acumulador no se modifica.
 
% ----------------------------------------------------------------------
\section{Pregunta 5: Franjas horizontales de colores en assembler Intel}
 
\textbf{Idea.} La pantalla de texto de $80 \times 25$ caracteres dispone de 16 colores (numerados de 0 a 15). La imagen del enunciado muestra 16 franjas, una por color: cada franja ocupa una fila y 16 columnas, y la primera es negra (por eso no se distingue del fondo). Se usan dos funciones de la interrupción \texttt{INT 10h} del BIOS:
 
\begin{itemize}
  \item \texttt{AH = 02h}: posiciona el cursor (\texttt{DH} = fila, \texttt{DL} = columna, \texttt{BH} = página).
  \item \texttt{AH = 09h}: escribe el carácter de \texttt{AL} con el atributo de \texttt{BL}, repetido \texttt{CX} veces desde la posición del cursor.
\end{itemize}
 
Cada franja se dibuja con el carácter de bloque lleno (\texttt{0DBh}) cuyo color de primer plano es el número de la fila. De este modo, la fila $n$ se pinta con el color $n$ y no hace falta una tabla de colores.
 
\begingroup
\footnotesize
\setlength{\tabcolsep}{5pt}
\begin{center}
\begin{tabular}{c c l | c c l}
  \hline
  \textbf{Fila} & \textbf{Color} & \textbf{Nombre} & \textbf{Fila} & \textbf{Color} & \textbf{Nombre} \\
  \hline
  0 & 0h & Negro        &  8 & 8h & Gris oscuro \\
  1 & 1h & Azul         &  9 & 9h & Azul claro \\
  2 & 2h & Verde        & 10 & Ah & Verde claro \\
  3 & 3h & Cian         & 11 & Bh & Cian claro \\
  4 & 4h & Rojo         & 12 & Ch & Rojo claro \\
  5 & 5h & Magenta      & 13 & Dh & Magenta claro \\
  6 & 6h & Marrón       & 14 & Eh & Amarillo \\
  7 & 7h & Gris claro   & 15 & Fh & Blanco \\
  \hline
\end{tabular}
\end{center}
\endgroup
 
\textbf{Programa} (\texttt{franjas.asm}):
 
\begin{figure}[H]
    \centering
    \includegraphics[width=0.86\textwidth]{img/franjasASM.png}
    \caption{Franjas.asm en EMU8086}
    \label{fig: franjas.asm}
\end{figure}
 
\textbf{Funcionamiento.}
 
\begin{itemize}
  \item Se fija el modo de texto 80$\times$25 con \texttt{AX = 0003h}, lo que también limpia la pantalla.
  \item \texttt{DH} lleva la fila y, a la vez, el número de color (0 a 15). En cada vuelta se posiciona el cursor en la columna 0 de esa fila.
  \item Se escriben 16 bloques llenos con el atributo \texttt{BL = DH}, es decir, una franja de 16 columnas del color de esa fila.
  \item Se incrementa \texttt{DH} y se repite hasta completar 16 franjas (filas 0 a 15).
  \item Se espera una tecla (\texttt{INT 16h}) para poder ver la pantalla y se termina con \texttt{INT 21h, AH = 4Ch}.
\end{itemize}
 
Se usa el color de primer plano porque funciona igual en cualquier entorno; los colores de fondo 8 a 15 dependen de si el modo de parpadeo del texto está activado.
 
\begin{figure}[H]
    \centering
    \includegraphics[width=0.86\textwidth]{img/franjasCOM.png}
    \caption{Ejecución de Franjas.asm en EMU8086}
    \label{fig: franjas.com}
\end{figure}
% ----------------------------------------------------------------------
\section{Pregunta 6: Manipulación de bits del registro CX}
 
Se pide, sobre el dato de CX (bits D7 a D0): invertir los bits impares D5, D3 y D1; borrar los bits pares D6, D4 y D2; y poner en 1 los bits D7 y D0. Cada acción se resuelve con una instrucción lógica y una máscara:
 
\begin{itemize}
  \item \textbf{Invertir} bits: \texttt{XOR} con 1 en los bits a invertir. Máscara $0010\ 1010_2 = 2Ah$ (D5, D3, D1).
  \item \textbf{Borrar} bits: \texttt{AND} con 0 en los bits a borrar y 1 en el resto. Máscara $1010\ 1011_2 = ABh$ (D6, D4, D2 en cero). Como CX tiene 16 bits, se usa $FFABh$ para no alterar el byte alto.
  \item \textbf{Poner en 1} bits: \texttt{OR} con 1 en los bits a poner. Máscara $1000\ 0001_2 = 81h$ (D7 y D0).
\end{itemize}
 
\begingroup
\footnotesize
\setlength{\tabcolsep}{3pt}
\begin{center}
\begin{tabular}{|c|c|c|c|c|c|c|c|c|}
  \hline
  \textbf{Bit} & D7 & D6 & D5 & D4 & D3 & D2 & D1 & D0 \\
  \hline
  \textbf{Acción} & poner 1 & borrar & invertir & borrar & invertir & borrar & invertir & poner 1 \\
  \hline
\end{tabular}
\end{center}
\endgroup
 
Las tres operaciones afectan a bits distintos, por lo que el orden en que se aplican no cambia el resultado.
 
\textbf{Programa} (\texttt{bits.asm}), con el dato de ejemplo del enunciado ($CX = 00B1h$, es decir, $1011\ 0001_2$ en D7 a D0):
 
\begin{figure}[H]
    \centering
    \includegraphics[width=0.86\textwidth]{img/bitsASM.png}
    \caption{bits.asm en EMU8086}
    \label{fig: bits.asm}
\end{figure}
 
Las cuatro primeras instrucciones son la solución; el bloque \texttt{imprimir} solo sirve para ver el resultado en DOSBox, donde no hay ventana de registros.
 
\textbf{Comprobación con el ejemplo:}
 
\begingroup
\small
\setlength{\tabcolsep}{4pt}
\begin{center}
\begin{tabular}{|l|l|c|}
  \hline
  \textbf{Paso} & \textbf{Efecto} & \textbf{Byte bajo de CX} \\
  \hline
  Dato inicial & --- & $1011\ 0001_2 = B1h$ \\
  \hline
  \texttt{XOR CX, 002Ah} & invierte D5, D3 y D1 & $1001\ 1011_2 = 9Bh$ \\
  \hline
  \texttt{AND CX, 0FFABh} & borra D6, D4 y D2 & $1000\ 1011_2 = 8Bh$ \\
  \hline
  \texttt{OR CX, 0081h} & pone en 1 D7 y D0 & $1000\ 1011_2 = 8Bh$ \\
  \hline
\end{tabular}
\end{center}
\endgroup
 
En este ejemplo el \texttt{OR} no cambia el valor porque D7 y D0 ya estaban en 1.
 
\textbf{Respuesta:} el resultado es $CX = 008Bh$, es decir, $1000\ 1011_2$ en D7 a D0. Para cualquier otro dato el byte alto no cambia; por ejemplo, $CX = 1234h$ queda en $128Bh$.
\begin{figure}[H]
    \centering
    \includegraphics[width=0.86\textwidth]{img/bitsCOM.png}
    \caption{Ejecución de bits.asm en EMU8086}
    \label{fig: bits.com}
\end{figure}

\end{document}